% TOP_PROBLEM: {"problemId": "problem.top500-omission-w045049-kobayashi-canonical-ampleness", "canonicalTitle": "Kobayashi Conjecture: Hyperbolic Compact Kähler Manifolds Have Ample Canonical Bundle", "releaseRank": 364, "primaryDomain": "geometry_topology", "primaryDomainLabel": "Geometry & topology", "displayStatus": "open_with_solved_subcases", "releaseStatus": "open_with_solved_subcases"}
% !TeX program = lualatex
% ATTEMPT_STATUS: unresolved
\documentclass[11pt]{article}
\usepackage[margin=25mm]{geometry}
\usepackage{fontspec}
\setmainfont{Latin Modern Roman}
\usepackage{amsmath,amssymb,mathtools}
\usepackage{unicode-math}
\setmathfont{Latin Modern Math}
\usepackage{xurl,hyperref}
\hypersetup{colorlinks=true,urlcolor=blue}
\setcounter{secnumdepth}{0}
\setlength{\parindent}{0pt}
\setlength{\parskip}{5pt}
\setlength{\emergencystretch}{3em}
\begin{document}
\section*{\texorpdfstring{Kobayashi Conjecture: Hyperbolic Compact Kähler Manifolds Have Ample Canonical Bundle}{Kobayashi Conjecture: Hyperbolic Compact Kähler Manifolds Have Ample Canonical Bundle}}\label{364}
\subsection{Definitions and mathematical statement}
A compact connected complex manifold is K\"ahler if it has a Hermitian metric whose associated real $(1,1)$-form is closed. The Kobayashi pseudodistance is the largest pseudodistance contracted by all holomorphic maps from the unit disk with its Poincar\'e metric. Hyperbolicity means it separates distinct points; in the compact setting this is equivalent to every holomorphic map ${\mathbb{C}}\to X$ being constant. The conjecture asserts
\[
 X\text{ compact K\"ahler and Kobayashi hyperbolic}
 \quad\Longrightarrow\quad K_X\text{ ample}.
\]
Here $K_X=\bigwedge^{\dim_{{\mathbb{C}}} X}\Omega_X^1$, and ample means that a positive power gives a projective embedding. Projectivity is thus part of the proposed conclusion, not an extra hypothesis. Negative curvature of a chosen metric is not assumed.
\par\smallskip\textit{Formulation and concept references:} \href{https://arxiv.org/pdf/2011.11379}{[S1]}.

\subsection{Short English statement}
Must a compact K\"ahler manifold with no nonconstant entire curves have an ample canonical bundle?
\subsection{Research attempt: separating hyperbolicity, bigness, and ampleness}
\paragraph{Scope and source audit.}
The hypothesis is compact K\"ahler hyperbolicity, with no projective embedding or negatively curved metric supplied. The primary exposition [S1], version 2, was checked on 27 September 2026. It distinguishes the projective nefness consequence of Mori theory, the stronger negative-holomorphic-sectional-curvature theorem, and the general conjecture. Its Lemma 5.1 treats the projective general-type case; its Section 4 explains why a negative-sectional-curvature metric cannot be assumed for every hyperbolic manifold. These are scope restrictions, not alternative formulations of the full target. The discussion below establishes several elementary subcases and reductions. Deep external inputs are identified where used, and the missing positivity step is retained.

In dimension zero the connected manifold is a point, and the conclusion holds. Henceforth let $n=\dim_{\mathbb C}X\ge1$. We use the compact Brody characterization stated in the problem: every entire map $\mathbb C\to X$ is constant. This characterization is not being asserted for arbitrary noncompact complex spaces.

\paragraph{1. What hyperbolicity immediately forbids.}
There is no nonconstant holomorphic map $\mathbb P^1\to X$. Indeed, restriction to the affine chart $\mathbb C\subset\mathbb P^1$ would be a nonconstant entire map: if it were constant on that open chart, the identity theorem would make the original map constant. Nor can there be a nonconstant map from an elliptic curve $E$ to $X$. Compose such a map with the surjective universal covering $\mathbb C\to E$.

More generally, if $C\subset X$ is an irreducible compact curve and $\widetilde C\to C$ its normalization, then
\[
 g(\widetilde C)\ge2.                                  \tag{1}
\]
The normalization map followed by inclusion is nonconstant, and compact Riemann surfaces of genus zero and one have the preceding forms. The assertion concerns geometric genus, so singular rational curves are excluded as well as smooth ones.

Also, every compact complex submanifold of $X$ is hyperbolic: an entire curve in it would compose to an entire curve in $X$. This elementary inheritance will constrain fibers later. It does not itself estimate the canonical degree of every subvariety. For a curve in a surface, the adjunction formula contains a self-intersection term, and in higher codimension the normal bundle contributes; genus information alone is not a formula for $K_X\cdot C$.

\paragraph{2. A complete proof in dimension one.}
Let $X=C$ be a compact connected Riemann surface of genus $g$. The preceding argument gives $g\ge2$, and
\[
 \deg K_C=2g-2>0.
\]
Here is an explicit embedding consequence rather than merely a numerical conclusion. A line bundle $A$ of degree at least $2g+1$ separates every length-two subscheme $Z\subset C$. By Serre duality,
\[
 H^1(C,A(-Z))^{\vee}=H^0(C,K_C\otimes A^{-1}(Z))=0,
\]
since the last line bundle has degree at most $-1$. The restriction map $H^0(C,A)\to H^0(Z,A|_Z)$ is therefore onto. It separates distinct points and tangent directions, and the usual very-ampleness criterion gives an embedding. Applying this to $A=K_C^3$ works because $6g-6\ge2g+1$ for every $g\ge2$. Thus $K_C$ is ample, proving the target in dimension one.

The converse hyperbolicity statement for such a curve follows from uniformization: its universal cover is the unit disk. An entire map lifts to a disk-valued entire map because $\mathbb C$ is simply connected, and Liouville's theorem makes the lift constant. Uniformization and the line-bundle embedding criterion are classical external background; the degree and duality argument above gives the claimed positive power explicitly.

\paragraph{3. Products and finite unramified quotients.}
Take compact curves $C_1,\ldots,C_n$ of genera at least two and put $Y=\prod_i C_i$. Every entire map to $Y$ has constant coordinate projections, so $Y$ is hyperbolic. The cotangent splitting gives
\[
 K_Y\simeq\bigotimes_i\operatorname{pr}_i^*K_{C_i}.      \tag{2}
\]
Choose one integer $a>0$ for which all $K_{C_i}^a$ are very ample; for instance $a=3$ works. The product of the resulting curve embeddings followed by the Segre embedding is defined by a subspace of $H^0(Y,K_Y^a)$ and embeds $Y$. Thus $K_Y$ is ample. This verifies arbitrarily high-dimensional examples without inferring positivity from an unconstructed curvature metric.

Now let $f:Y\to X$ be a finite \emph{\'etale} surjective map of connected compact complex manifolds, with $X$ K\"ahler. Analytically $f$ is a covering map. An entire curve in $X$ lifts to $Y$, so if $Y$ is hyperbolic then $X$ is hyperbolic. Conversely, if $X$ is hyperbolic, composition shows that an entire curve in $Y$ has image in one finite fiber, hence is constant. Thus hyperbolicity is equivalent on the two spaces in this setting.

The absence of ramification gives $K_Y=f^*K_X$. If $K_Y$ is ample, the standard finite-descent theorem for ampleness implies that $K_X$ is ample. One can see the numerical content after projectivity is established by that theorem: for any positive-dimensional integral $V\subset X$ and an irreducible component $W$ above it,
\[
 (f^*K_X)^{\dim V}\cdot W
   =\deg(W/V)\,K_X^{\dim V}\cdot V>0,
\]
and Nakai--Moishezon yields ampleness downstairs. Finite descent, including projectivity in the analytic setting for a finite quotient of a projective manifold, is an external standard result here; the intersection computation is not a substitute for it on an initially nonprojective space.

In particular, a free finite group quotient of a product of the curves above satisfies the conjecture. Dropping ``\'etale'' changes both arguments: an entire curve need not lift across branch points, and the canonical bundle formula acquires a ramification divisor. The unramified hypothesis is doing real work.

\paragraph{4. The projective nefness reduction.}
Suppose now that $X$ is smooth and projective. We use the following consequence of the cone theorem as external input: if $K_X$ is not nef, $X$ contains a rational curve on which $K_X$ has negative degree. Paragraph 1 excludes such a curve. Therefore
\[
 X\text{ projective and hyperbolic}\quad\Longrightarrow
 \quad K_X\text{ nef}.                                \tag{3}
\]
This is a conclusion about nonnegative degrees on curves. It is not ampleness. A nef line bundle may have zero top self-intersection or contract curves. The trivial line bundle on a positive-dimensional projective variety already shows that nefness alone implies neither bigness nor an embedding.

For a nef line bundle on a projective $n$-fold, bigness is equivalent to positivity of the top self-intersection. Thus a natural sufficient next step in the projective case is
\[
 \int_X c_1(K_X)^n>0.                                  \tag{4}
\]
Given (4), the standard general-type lemma in [S1, Lemma 5.1] says that a smooth projective variety with no rational curves has ample canonical bundle. That lemma uses the basepoint-free and log cone theorems; it is not proved merely by testing $K_X$ on the rational curves that are already absent. We use it only as an explicitly external reduction.

For clarity, the numerical obstruction in its proof is this. If the canonical morphism contracts a curve $C$, then $K_XC=0$. Bigness supplies an effective rational divisor numerically equivalent to $K_X-\varepsilon H$ for a small positive $\varepsilon$ and an ample $H$, so this divisor has negative intersection with $C$. A sufficiently small boundary then allows the log cone theorem to produce a rational curve, contradiction. Without bigness the effective divisor with that class is unavailable.

The same discussion cannot simply begin on a compact nonprojective K\"ahler manifold by declaring the projective cone theorem applicable. The problem explicitly asks us to obtain projectivity. A route through a big canonical line bundle would produce a Moishezon manifold; the classical theorem that K\"ahler plus Moishezon implies projective would then permit the projective argument. Obtaining that bigness from hyperbolicity remains a missing input.

\paragraph{5. An elementary surface calculation once bigness is supplied.}
There is a shorter proof of the last ampleness step for a smooth projective hyperbolic surface with $K_X$ big. By (3), $K_X$ is nef, and hence $K_X^2>0$. Suppose that an irreducible curve $C$ has $K_XC=0$. The Hodge index theorem says that the intersection pairing is negative definite on the orthogonal complement of a class of positive square. Since $C$ is numerically nonzero (it has positive intersection with an ample divisor), it follows that $C^2<0$.

Adjunction gives
\[
 2p_a(C)-2=C^2+K_XC=C^2.                               \tag{5}
\]
The arithmetic genus of an integral projective curve is nonnegative. Consequently (5) forces $p_a(C)=0$ and $C^2=-2$. The normalization formula $p_a(C)=g(\widetilde C)+\sum_p\delta_p$ then says that $C$ is a smooth rational curve, a contradiction. We have proved $K_XC>0$ for every irreducible curve and $K_X^2>0$. Nakai--Moishezon now makes $K_X$ ample.

This calculation checks the singular-curve case explicitly: it would be incorrect to insert the geometric genus in adjunction before accounting for the singularity corrections. It also displays precisely where the assumed bigness enters, through the positive square needed for Hodge index. The argument does not show that a hyperbolic surface has big $K_X$ without further classification or analytic input, and does not claim a new proof of all known surface results.

\paragraph{6. What an assumed canonical fibration would leave to do.}
Work again projectively, and suppose in addition that $K_X$ is semiample. A sufficiently divisible positive power defines a morphism with connected fibers
\[
 f:X\longrightarrow Y,\qquad K_X^a=f^*A,               \tag{6}
\]
where $A$ is ample on the normal image and the connected-fiber assertion is obtained by Stein factorization. If $\dim Y=n$, the canonical bundle is big, and the preceding external general-type lemma applies. If $\dim Y<n$, consider a smooth general fiber $F$ over the smooth locus where $f$ is smooth. Generic smoothness in characteristic zero supplies such fibers.

The normal bundle of $F$ is trivial, being the pullback of the tangent space of the base at its point. Adjunction therefore gives $K_F=K_X|_F$. Restricting (6) yields
\[
 K_F^a\simeq\mathcal O_F.                              \tag{7}
\]
On the other hand $F$ is a compact hyperbolic manifold by inheritance. Thus this proposed route would need to rule out positive-dimensional smooth hyperbolic fibers with torsion canonical bundle.

For a one-dimensional fiber this is immediate: (7) forces genus one, contrary to paragraph 1. For a complex torus of any positive dimension it is also immediate: a nonzero complex-linear map $\mathbb C\to\mathbb C^q$ composed with the quotient by its lattice is an entire nonconstant curve. The conclusion remains true for a finite \emph{\'etale} quotient of such a torus by the same covering argument.

These observations do not classify all varieties with torsion canonical bundle. In particular, (7) does not say that $F$ is a torus. Moreover, the semiampleness assumed in (6) is itself a substantial extra input in arbitrary dimension. We therefore obtain a conditional fibration obstruction, not a deduction of abundance or of the original conjecture.

\paragraph{7. A curvature identity that proves a stronger-hypothesis case.}
Suppose $X$ carries a K\"ahler metric with negative definite Ricci form. In local coordinates let $(g_{i\bar j})$ be its positive Hermitian matrix. The induced metric on the canonical frame $dz_1\wedge\cdots\wedge dz_n$ has squared norm $(\det g)^{-1}$. With the convention $i\Theta(L,h)=-i\partial\bar\partial\log h$, we obtain
\[
 i\Theta(K_X)=i\partial\bar\partial\log\det g
             =-\operatorname{Ric}(g).                 \tag{8}
\]
Hence $K_X$ admits a positive Hermitian curvature form. Kodaira's embedding theorem implies that it is ample and that $X$ is projective. Formula (8) fixes the sign: negative Ricci corresponds to positive canonical curvature, not to positivity of the anticanonical bundle.

This is an entirely sufficient extra hypothesis, but hyperbolicity does not provide the metric in (8). The Wu--Yau and K\"ahler extensions recorded in [S1] obtain canonical ampleness under negative holomorphic sectional curvature through additional analysis. Replacing the given hyperbolicity assumption by that curvature assumption would change the problem. Nor is there a legitimate pointwise argument that every negative holomorphic sectional curvature metric already has the desired Ricci sign; the curvature contractions and their sign implications require care.

\paragraph{8. Stress test: ample canonical bundle is not enough for hyperbolicity.}
Consider the explicit quintic surface $Z\subset\mathbb P^3$ with equation
\[
 F=x_0^4x_2+x_2^5+x_1^4x_3+x_3^5=0.                   \tag{9}
\]
It is smooth. Indeed the first and third partial derivatives are $4x_0^3x_2$ and $x_0^4+5x_2^4$. Their simultaneous vanishing forces $x_0=x_2=0$: if either factor of the first is zero, the second forces the other coordinate to vanish. The other two partial derivatives similarly force $x_1=x_3=0$. Thus the gradient has no zero in projective space.

Adjunction gives $K_Z=\mathcal O_Z(1)$, which is ample. Nevertheless the line $x_2=x_3=0$ lies in $Z$, so $Z$ is not hyperbolic. This is not a counterexample to the conjecture, whose implication goes the other way. It prevents one from treating canonical positivity and hyperbolicity as equivalent and then arguing circularly through a supposed converse. The saved symbolic diagnostic checks the partial derivatives and the restriction to the line; smoothness follows from the exact argument above, rather than from numerical sampling.

\paragraph{9. First gap and outcome.}
The arguments prove the target for curves, for products of hyperbolic curves, and for their finite free quotients; they prove the final positivity step for projective hyperbolic manifolds once bigness is supplied, with an explicit surface argument. They also identify a sufficient Ricci-curvature hypothesis and a conditional restriction on a semiample canonical fibration.

For a general projective hyperbolic manifold, the first missing step in this route is to force canonical bigness, or an alternative condition strong enough to upgrade (3) to ampleness. For the full K\"ahler formulation, projectivity must additionally be earned rather than assumed. Neither the absence of rational and elliptic curves nor the degree computations on known examples supplies these missing statements.

\textbf{Outcome: UNRESOLVED.} The reductions retain both central quantifiers: every compact K\"ahler hyperbolic manifold, and existence of an ample canonical bundle without an initially chosen negative-curvature metric. The partial constructions and obstructions do not settle that general implication.

\subsection{Sources}
\begin{itemize}
\item[S1]S. Diverio, Kobayashi hyperbolicity, negativity of the curvature and positivity of the canonical bundle, version 2 (2020).
\url{https://arxiv.org/pdf/2011.11379}.
\textit{Formulation source.}
\end{itemize}
\end{document}
